\documentclass[sigconf,screen]{acmart}
\pdfsuppressptexinfo=-1
\AtBeginDocument{
  \providecommand\BibTeX{{\normalfont B\kern-0.5em{\scshape i\kern-0.25em b}\kern-0.8em\TeX}}}
\usepackage{booktabs}
\usepackage{amsmath}
\setcopyright{cc}
\setcctype{by}
\acmDOI{10.1145/3843282.3844447}
\acmYear{2026}
\copyrightyear{2026}
\acmISBN{979-8-4007-2985-0/2026/10}
\acmConference[AgenticDev '26]{Proceedings of the 1st International Workshop on Agentic AI for Next-Generation Software Development}{October 12--16, 2026}{Munich, Germany}
\acmBooktitle{Proceedings of the 1st International Workshop on Agentic AI for Next-Generation Software Development (AgenticDev '26), October 12--16, 2026, Munich, Germany}
\acmSubmissionID{asews26agenticdevmain-p195-p}
\received{2026-07-25}
\received[accepted]{2026-08-20}

\begin{document}

\title{They Do Not Fire Late, They Barely Fire}
\subtitle{Four Surface-Behavior Stopping Baselines for Long-Horizon Coding Agents,
and How to Report One}
\author{Simarjot Khanna}
\affiliation{
  \institution{Independent Researcher}
  \country{Canada}
}
\email{simarkhanna76@gmail.com}

\begin{abstract}
Long-horizon coding agents frequently waste budget by acting without progressing. A growing family
of stop-signals proposes to abort such runs early, but each is evaluated on its own corpus against
its own baseline, leaving their limits unknown. We operationalize four stop-signal families from
published designs as inexpensive fixed-threshold detectors and evaluate them by replaying public
agent trajectories; every claim is scoped to these author-built operationalizations, not to the
published systems themselves. A retrospective Oracle Stop, the last step at which a run moved
closer to the benchmark's reference patch, grounds the timing and waste analyses; the headline
results need only run outcomes and detector fire steps. Held to a strict false-abort budget of
5\%, these detectors fail not by firing late but by remaining silent on the large majority of
doomed runs. The diagnosis replicated on four preregistered held-out configurations, and a
threshold calibrated on one deployment either fell silent or violated the budget on others. The
mechanism is that doomed and recoverable runs produce largely overlapping alarm-score
distributions, and their separability stops improving by mid-run, so waiting does not sharpen the
signal; in a difficulty-matched sample, task difficulty does not produce the separation. From this
we derive a falsifiable separation target for future detectors, and a three-part reporting
protocol: coverage at a stated false-abort budget, run-level separation, and whether the operating
threshold survives a change of deployment.
\end{abstract}

\begin{CCSXML}
<ccs2012>
<concept><concept_id>10011007.10011074.10011099</concept_id>
<concept_desc>Software and its engineering~Software verification and validation</concept_desc>
<concept_significance>500</concept_significance></concept>
<concept><concept_id>10010147.10010178</concept_id>
<concept_desc>Computing methodologies~Artificial intelligence</concept_desc>
<concept_significance>300</concept_significance></concept>
</ccs2012>
\end{CCSXML}
\ccsdesc[500]{Software and its engineering~Software verification and validation}
\ccsdesc[300]{Computing methodologies~Artificial intelligence}

\keywords{coding agents, early termination, stop-signals, agent liveness,
empirical software engineering}

\maketitle

\section{Introduction}
An autonomous coding agent is given a budget and a task and decides for itself when it is done.
When the task is out of reach it frequently does something worse than fail: it keeps going,
re-reading files it has read, re-applying edits it has tried, rerunning tests it has watched fail,
consuming budget while producing nothing that moves it toward a solution. The run is not deadlocked;
it simply is not converging. This is now widely reported \cite{bouzenia2025understanding,
understandingcodeagent2025, coherencecollapse2026, resolveratehides2026}, and a parallel line of
work proposes to cut such runs short, each reporting savings on its own corpus against its own
baseline \cite{bagen2026, earlydiagnosis2026, agentstop2026, shp2026, eet2026}.

What the field cannot currently do is say why those proposals stop where they do. A stop-signal is
a predictor of a single latent event, \emph{the agent will not recover from here}, and predictors are
diagnosed against ground truth. There has been none. Whether the agent was still making progress at
step $t$ has been answerable only by an LLM judge \cite{shp2026} or a heuristic warning
\cite{earlydiagnosis2026}, which are themselves instances of the fallible signal under study.

For coding agents a retrospective progress proxy is available and free. A benchmark instance ships
with a reference (gold) solution; at each step we measure how close the agent's accumulated edits
have come to it and read off the last step at which they came closer than ever before. We call that
step the \textbf{Oracle Stop} $t^{*}$. It is not a deployable stopping rule, because it uses the
answer; it is a measuring instrument, mirroring the oracle upper bounds used to study single-model
reasoning-token efficiency \cite{beyondaccuracy2026}, lifted to multi-step agents. Grounded on it we
place detector families operationalized from published stop-signals on a common
savings-versus-false-abort frontier, by pure retrospective replay of public trajectories. The
oracle grounds the timing and waste analyses; the coverage and transfer results ahead are computed
from outcome labels and detector fire steps alone (Section~\ref{sec:oracle}).\footnote{Artifact
(harness, frozen specification, recorded results, verifier):
\url{https://github.com/s1mar/stop-signal-baselines}}

\paragraph{The diagnosis.} These detectors are silent more often than they are slow: at the strict operating points a
deployment would want, they almost never fire at all. On our development corpus of 708
long-horizon trajectories, holding false aborts on resolved runs to 5\%, the best family recovers
0.129 of failed-run compute, fires at all on 0.18 of failed runs, and has an active alarm at $t^{*}$
on under 0.03 of them (the last of these, unlike the first two, depends on the oracle). Relax the
constraint entirely and the three mechanical families cover 0.99 to 1.00. What binds is neither the time budget nor the
stagnation statistic, but the ability to tell a doomed run from a recoverable one that looks the
same. Everything above is a claim about the four inexpensive fixed-threshold baselines we
operationalize, and about nothing else: not about the published systems whose signal types they
draw on, and not about stop-signals as a class.

\paragraph{The evidence that it generalizes.} That diagnosis was discovered on one corpus during
adversarial review, so we treated it as exploratory, wrote a frozen analysis
specification with a stated decision rule, and tested it once on four held-out configurations
spanning three model families, two scaffolds and three task sets. It replicates on 4 of 4: median
coverage at the 5\% budget across the 16 held-out detector-configuration pairs is 0.043 against 0.16
to 0.18 for the mechanical families on the development corpus, and 15 of the 16 pairs are lower off-corpus than the same family
on the development corpus. A
separate preregistered check shows that operating thresholds fitted on one configuration do not
transfer: applied unchanged, they either fire on almost nothing or exceed the false-abort budget,
in the worst case by nearly thirteen times.

\paragraph{The mechanism, and why it is not an artifact of the operating point.} Coverage and
latency are both consequences of thresholding one score, so a purely threshold-based argument would
be circular. We therefore measure the underlying quantity threshold-free: the run-level alarm-score
distributions of failed and resolved runs overlap by 0.75 to 0.92 and separate at only 0.55 to 0.62
AUC (0.60 to 0.62 for the three mechanical families, 0.546 for self-report), and that separability
peaks by mid-run and then stops improving. Waiting does not sharpen the signal; it only spends
budget.

\paragraph{From a negative result to a design target.} Overlap is a quantity a future detector can
be measured against. Section~\ref{sec:target} converts it into a falsifiable requirement: covering
half of doomed runs at a 5\% false-abort budget needs run-level separation of 0.91 to 0.95 AUC on
the empirical score shapes, and 0.88 under the more forgiving equal-variance Gaussian
idealization, against 0.617 measured here. That is a large improvement, well beyond what threshold tuning reaches.

\paragraph{Contributions.}
\begin{itemize}
\item \textbf{A reporting protocol for stopping rules}, which is the transferable output of the
study: coverage at a stated false-abort budget, run-level separation with an instance-clustered
interval, and whether the operating threshold survives a change of deployment. Each of the three is
shown here to be load-bearing, none needs more than logged detector scores, resolved/failed outcome
labels and task identifiers, and none is standard practice.
\item A diagnosis, replicated on four preregistered held-out configurations, that the
surface-behavior stop-signals we build fail by coverage collapse rather than by latency, and that
their thresholds do not transfer across deployments.
\item A threshold-free mechanism for that collapse: run-level score overlap of 0.75 to 0.92 and
separability that plateaus by mid-run, which rules out the operating point as the explanation and
shows that waiting does not help, and a difficulty-matched corpus showing the separation is not an
artifact of which tasks are hard.
\item A falsifiable design target derived from the mechanism, stating the run-level separation a
usable stop-signal of this class would need at a given false-abort budget, for score distributions
shaped like the ones we observe.
\item A measurement defense: the cumulative gold-patch oracle is validated against a gold-blind
three-vendor panel and against a more ``principled'' current-state alternative that degenerates,
and we separate the claims that depend on the oracle from those that do not.
\end{itemize}

\section{Related Work}
\paragraph{Agent failure, stopping, and wasted computation.}
A fast-growing body of work dissects why coding agents fail: \citet{bouzenia2025understanding} link
failure to thought-action misalignment and repetitive action cycles,
\citet{understandingcodeagent2025} contrast success and failure trajectories,
\citet{coherencecollapse2026} identify agents that reach a correct fix and then degrade it, and
\citet{resolveratehides2026} argue resolve rate hides trajectory structure. A parallel line
proposes policies to cut such runs short: budget-aware early stopping is reported to save 28 to
64\% of tokens on failed trajectories \cite{bagen2026}, a large fraction of tokens
is spent after a first failure warning \cite{earlydiagnosis2026}, and others terminate on token
log-probabilities \cite{agentstop2026}, on semantic convergence between successive drafts,
optionally gated by a judged quality score \cite{shp2026}, or on learned experience \cite{eet2026}. The first group diagnoses the phenomenon
whose cost we quantify; the second optimizes policies without a shared reference against which
their limits can be characterized. We do not propose a better policy, and we do not evaluate those
systems. We build cheap operationalizations of the same signal types and ask what limits them,
answering with a mechanism and a target.

\paragraph{Reference-grounded analysis.}
Grounding trajectory analysis in the reference patch is not itself novel:
\citet{coherencecollapse2026} use it and describe the paradigm as concurrent work.
\citet{beyondaccuracy2026} define an oracle early-stop upper bound from the gold answer for
single-model reasoning chains. We do not claim gold grounding as a contribution. We claim the
diagnosis built on top of it, its preregistered replication, and the honest characterization of the
oracle's error.

\paragraph{Liveness.}
An agent that acts forever without progressing violates a classical liveness property
\cite{lamport1977proving, alpern1985defining}; recent work names this \emph{semantic livelock} and
proposes an embedding-diversity monitor for it \cite{zombieagents2026}. Our results locate that
line's difficulty: the monitor sees stagnation, and stagnation is not what separates a doomed run
from a recoverable one.

\section{The Instrument}
\paragraph{Setup.} A trajectory is a sequence of agent steps $s_1,\dots,s_n$. A subset of steps are
edit actions that write content to files. Each benchmark instance has a gold patch $G$, a unified
diff whose added lines define a target file set $F_G$ and a target token set $T_G$ (the code
identifiers added by $G$).

\paragraph{Progress.} At edit step $t$ let $A(t)$ be the set of code identifiers the agent has
written into the gold target files $F_G$ up to $t$. Restricting to $F_G$ discards throwaway
scaffolding (reproduction scripts, debug files) that would otherwise dominate the signal. We define
primary progress as token recall
\begin{equation}
p(t) = |A(t) \cap T_G| \,/\, |T_G| \in [0,1],
\end{equation}
which is robust to reformatting and to the textual variation between two correct fixes, unlike
exact line or full-patch matching. As a secondary, semantic view we also compute
$p_2(t)=\cos(\phi(A_t),\phi(G))$, where $\phi$ is a frozen sentence encoder
\cite{reimers2019sentencebert} and $A_t$ and $G$ are the text the agent and the reference patch have
added. When the two views disagree on the Oracle Stop by more than
20\% of the trajectory length we mark the trajectory \emph{contested} and report the rate rather
than resolving it.

\paragraph{Oracle Stop.} The Oracle Stop is the last edit step that reached a new maximum progress:
\begin{equation}
t^{*} = \max\{\, t : p(t) > p(t') \ \forall\, t' < t \,\}.
\end{equation}
with $t^{*}$ pinned at the run's first edit step when no step attains positive progress (the
zero-overlap stratum of Section~\ref{sec:oracle}). On a failed run, every step after $t^{*}$ produced no further movement
toward the known solution and no resolution. Because $p$ is cumulative, $t^{*}$ is a \emph{last-improvement} point, not the
onset of a decline; Section~\ref{sec:oracle} measures how often failed runs decline at all, and
tests this instrument against a current-state alternative that can.

\paragraph{Detector quantities.} A detector $d$ fires at the first step whose alarm score crosses a
threshold $\theta$. Sweeping $\theta$ traces a curve in three quantities:
\begin{itemize}
\item \textbf{false-abort rate}, the fraction of \emph{resolved} trajectories the detector would
have aborted before completion; this is the deployment risk and we hold it to a budget;
\item \textbf{coverage}, the fraction of failed runs on which the detector fires at all;
\item \textbf{saved fraction}, the compute recovered by aborting at the fire point, as a fraction
of total failed compute.
\end{itemize}
We also report \textbf{Oracle Regret} $R_d = t_d - t^{*}$, how many steps after the Oracle Stop the
detector fires. Note which of these use the oracle. False-abort rate, coverage and the saved
fraction are computed from the run's outcome label and the step at which the detector fires; none
of them reads $t^{*}$. Only the Oracle Regret, the reference level defined next, and the ``alarm
active at $t^{*}$'' statistic do. This separation matters for the argument and we return to it in
Section~\ref{sec:oracle}.

\paragraph{The reference level.} For a failed trajectory of $n$ steps, the post-peak region is
$n-1-t^{*}$ steps, and the corpus-level \textbf{Post-Peak Ratio} is the total post-peak compute over
total failed compute, $\mathrm{PPR} = \sum_i (n_i - 1 - t^{*}_i) / \sum_i n_i$. PPR is the saved
fraction of a detector that fires exactly at $t^{*}$ on every failed run and never on a resolved
one. We use it only as a reference level on the frontier plot. It measures where a failed run's
compute falls relative to its last measured peak and nothing more, so it is neither a waste
magnitude nor an upper bound: a detector firing before $t^{*}$ saves strictly more. We report it, with its
mixture structure, in Section~\ref{sec:collapse}, and we do not build any claim of this paper on
its magnitude.

\section{Study Design}
\label{sec:design}
\paragraph{Corpus.} We stream the public Nebius \textsc{SWE-agent-trajectories} corpus
\cite{nebius2024trajectories} (80{,}036 runs of the SWE-agent scaffold \cite{yang2024sweagent} over
GitHub issues, three action-generator models) and join each trajectory to its gold patch from
\textsc{SWE-bench-extra}, the extended instance set released with that corpus
\cite{nebius2024trajectories}, whose task format follows SWE-bench \cite{jimenez2024swebench}. We
take a bounded, class-balanced sample, and the reduction to 708 is a
scope-and-joinability funnel: scanning the first 24{,}000 stream rows, we keep trajectories with at
least 15 agent steps, up to 450 failed and 300 resolved, of which 746 join to a gold patch (446
failed, 300 resolved). Two groups are then excluded because the progress instrument cannot score
them: 9 resolved runs whose gold patch adds no identifier tokens, and 29 failed runs in which no
edit action could be parsed. This leaves 708 trajectories (417 failed, 291 resolved, 80 unique
task instances; 45 instances contribute a failed run, 42 a resolved run and 7 both). The 15-step
floor is a scope choice: it
selects the long-horizon regime this paper is about and excludes short runs where stopping matters
less. All claims here are scoped to that regime. No agent is run: the frontier and every mechanical
detector are pure retrospective replay and reproduce on a laptop. Only the optional LLM-judge rows
and the gold-blind validation panel call inference.

\paragraph{Detector families.} We implement five evaluable stop-signal families behind one
interface. These are reimplementations \emph{in spirit}, not the published systems, and one
is an explicit stand-in, so what the study compares is detector families as we operationalize them,
not the original artifacts. Conclusions about stop-signals should be read at that grain.
(1)~\emph{Syntactic repetition}: recurrence of the normalized action string in a window. Repeated
action cycles are a documented failure mode \cite{bouzenia2025understanding} and self-reflective
agents such as Reflexion \cite{shinn2023reflexion} are designed to break out of them; the counter
is our own minimal detector for that behavior, not a reimplementation of either.
(2)~\emph{Embedding diversity}:
the Convergence Monitor \cite{zombieagents2026}, mean pairwise cosine similarity of the last $w{=}5$
action embeddings from a frozen encoder. (3)~\emph{Process-reward proxy}: the length of the current
run of steps whose observation shows a failure signal and no pass signal, a transparent stand-in for
a trained process reward model \cite{lightman2023verify}. (4)~\emph{Self-report}: accumulation of
give-up and uncertainty markers in the agent's own text, in the spirit of budget-aware
self-assessment \cite{bagen2026}. (5)~\emph{LLM-as-judge}: a small local model reads the issue and
the last few action-observation pairs and labels the step progressing or stuck \cite{shp2026};
because it needs per-step inference it is evaluated on a matched subset.
(6)~\emph{Log-probability confidence} \cite{agentstop2026} cannot be evaluated by replay: token
log-probabilities are not present in the public corpus. We report it as unavailable instead of
substituting a proxy.

\paragraph{Exploratory and confirmatory split.} The coverage-collapse diagnosis was found on the
development corpus during adversarial review, on the same data used to shape metric definitions and
controls. We label it \textbf{exploratory}. Before computing any held-out result we froze an
analysis specification: the hypotheses (coverage is low at a strict budget and rises as the
constraint relaxes; the dominant loss is coverage rather than latency), all four held-out
configurations with no subset selection, and unchanged detectors and threshold grids. The
specification also fixed the full list of quantities to report for every configuration, a
budget-sensitivity analysis to rule out dependence on the 5\% point specifically, the
threshold-transfer test, and a decision rule under which the headline would have been withdrawn or
narrowed. Section~\ref{sec:transfer} is the
\textbf{confirmatory} test. We report every configuration, including unfavorable rows.

\paragraph{How to read the measurements.} A deployable detector must hold the false-abort rate to
a small budget, because aborting runs that would have succeeded is the cost a deployment cannot
pay. Under that constraint, good performance is high coverage and a high saved fraction; a
detector can fail by firing late on the runs it covers, or by not firing at all. The run-level AUC
and overlap coefficient read the same signal without any threshold: separation near chance means
doomed and recoverable runs receive nearly the same scores, so no threshold can cover many doomed
runs within the budget.

\section{The Coverage Collapse}
\label{sec:collapse}

\begin{figure}[t]
\centering
\includegraphics[width=0.90\columnwidth]{figures/frontier.pdf}
\caption{Detector families against the Oracle Stop on the development corpus. Each curve sweeps a
detector's threshold; the dashed line is the PPR reference level. In the shaded deployable region
every family sits near the floor. The curves cross, and with instance-clustered intervals no
ordering among the three mechanical families is resolved.}
\Description{Line plot of saved fraction of failed-run compute against false-abort rate on
resolved runs, one curve per detector family. All curves stay near zero saved fraction inside the
shaded five percent false-abort region and only rise once the false-abort rate is large.}
\label{fig:frontier}
\end{figure}

\begin{table}[t]
\centering
\small
\setlength{\tabcolsep}{3pt}
\caption{Development corpus (417 failed, 291 resolved, 80 task instances). Saved fraction of failed
compute at three false-abort budgets, with coverage and median Oracle Regret at the 5\% budget.
Intervals are bootstrap 95\% CIs clustered by task instance with the operating threshold reselected
inside each replicate; for syntactic repetition and embedding diversity the interval reaches zero,
because on some resamples no threshold meets the budget. No detector is separated from any other, so we mark no winner.}
\begin{tabular}{lcccccc}
\toprule
& \multicolumn{3}{c}{Saved at false-abort} & \multicolumn{3}{c}{At the 5\% budget} \\
\cmidrule(lr){2-4}\cmidrule(lr){5-7}
Family & $\le$5\% & $\le$10\% & $\le$20\% & 95\% CI & Cov. & Regret \\
\midrule
Syntactic repetition & 0.129 & 0.160 & 0.160 & 0.000--0.191 & 0.18 & 10.0 \\
Embedding diversity  & 0.117 & 0.193 & 0.272 & 0.000--0.205 & 0.16 & 10.5 \\
Process-reward proxy & 0.125 & 0.168 & 0.218 & 0.070--0.193 & 0.18 & 11.0 \\
Self-report          & 0.049 & 0.101 & 0.101 & 0.026--0.095 & 0.07 & 12.0 \\
\midrule
LLM-as-judge & \multicolumn{6}{c}{\emph{matched subset only; see text}} \\
Log-prob confidence & \multicolumn{6}{c}{\emph{unavailable in this corpus}} \\
\bottomrule
\end{tabular}
\label{tab:frontier}
\end{table}

Table~\ref{tab:frontier} and Figure~\ref{fig:frontier} place the evaluable families on the
frontier. At the 5\% false-abort budget the best of them recovers 0.129 of failed-run compute; a
permissive 20\% budget reaches 0.272. With the bootstrap clustered by task instance and the
threshold reselected inside each replicate, no family is separated from any other, and for
syntactic repetition and embedding diversity the interval reaches zero, because on some
resamples no threshold meets the budget at all.
We report the level, which is robust, and decline to rank the detectors, which is not.

\paragraph{The loss is non-firing.} The natural reading of a ten-step regret is that detectors are
late. The compute goes elsewhere. At the 5\% operating point the best family fires
\emph{at all} on 0.18 of failed runs, and its alarm is active at $t^{*}$ on under 0.03 of them,
rising to about 0.20 ten steps later. A detector silent on four failed runs in five cannot recover
much compute at this budget, and no choice of threshold that respects the budget changes that. The
signal itself is present, though it does not recover uniformly: as the budget relaxes, the best
coverage rises to 0.24 at 10\% and 0.33 at 20\%, and with the constraint removed entirely the three
mechanical families cover 0.99 to 1.00 of failed runs and save 0.88 to 0.97 of their compute.
Syntactic repetition, however, sits flat at 0.23 across the 10\% and 20\% budgets. Self-report is the exception and is genuinely blind, never
exceeding 0.34 coverage at any budget: agents seldom announce that they are stuck.

This comparison is suggestive rather than decisive on its own, because coverage and latency are two
readings of one thresholded score, and any thresholded score can be forced to alarm everywhere.
Section~\ref{sec:mechanism} therefore measures the underlying separation without reference to any
threshold.

\paragraph{The collapse survives a more forgiving false-abort rule.} Counting every alarm before
completion on a resolved run as a lost resolution is the safe deployment assumption, but it is
pessimistic: a resolved run that has already written the patch which makes the tests pass, and is
then running tests or cleaning up, loses little if aborted. Because the operating threshold is set
by the extreme tail of the resolved population, a few such cases could in principle move it and with
it the whole frontier. We therefore recompute every operating point under two more forgiving
accountings, both by replay and neither using the gold patch: \emph{lenient}, which exempts alarms
at or after a resolved run's last applied edit, and \emph{edit window}, which starts the exemption
three steps earlier still. All three accountings sweep the same threshold grid as
Table~\ref{tab:frontier}, so the strict column reproduces it exactly and the comparison is
like-for-like. Relaxing the rule does buy something, but less than it would need to. Under the
lenient rule the process-reward proxy rises from 0.125 to 0.144 saved and embedding diversity from
0.117 to 0.142, while syntactic repetition and self-report are unmoved. Under the deliberately
over-generous edit window the process-reward proxy reaches 0.168 saved at 0.235 coverage and
syntactic repetition 0.160 at 0.225. Taking the best family under each accounting, the most
generous convention we can justify lifts saved compute from 0.129 to 0.168 and coverage from 0.18
to 0.24. The collapse is therefore a
property of the signal and not of the accounting convention, and the numbers in
Table~\ref{tab:frontier} are, as we assumed, a lower bound on what these families could achieve.

\paragraph{Reference level and its mixture.} For completeness we report the Post-Peak Ratio, which
sets the dashed line in Figure~\ref{fig:frontier}. It is 0.754 over all 417 failed runs (95\% CI
0.703 to 0.796, clustered by instance; 0.694 as a character proxy), but it is a mixture and should
not be read as one phenomenon. Of the 417, 205 never achieve any token overlap with the reference
patch: their Oracle Stop sits at the trajectory start, their PPR is 0.906, and they carry 0.571 of
all failed compute, so this never-on-track stratum contributes the larger share of post-peak
compute. The other 212 reach the patch, peak near the midpoint (median $t^{*}/n$ 0.511),
and have a PPR of 0.552 (95\% CI 0.480 to 0.622). Only the second population is stalling after a
near-approach.

\paragraph{The collapse is not an artifact of the mixture.} Recomputing the frontier separately on
each stratum, with the resolved set held fixed so false-abort keeps one denominator, barely moves
detector performance: at the 5\% budget syntactic repetition saves 0.118 on reached-reference runs
and 0.138 on the rest, and the other families differ by under 0.02 between strata. Compared within
population, 0.118 against that stratum's own 0.552 reference level is 0.21 of its post-peak
region, so about four fifths of that stratum's post-peak compute still goes unrecovered.
The detectors are not merely picking off runs that were never on track, and the diagnosis does not
depend on which population a reader treats as the legitimate target. The split itself reads the
reference patch, so this is a robustness check on the oracle-free result and not a second one.

\paragraph{The LLM judge is not resolved either way.} Because a judge needs per-step inference it is
evaluated on a matched subset (186 failed, 123 resolved) with every family recomputed on the same
trajectories. The two judges are open-weight coding models served locally through Ollama under an
identical prompt: a 30B Mixture-of-Experts model with 3B active (Apache-2.0) and a 33B MoE with 3B
active (OpenMDW-1.1).\footnote{Exact model tags, weight digests and the judge prompt are in the
artifact.} The conservative of the two, which
calls 18\% of steps stuck, attains the highest point estimate at the 5\% budget (0.159 saved, median
regret 7.0 steps against 9.5 to 13 for the mechanical families); the eager one, at 32\% stuck, cannot
satisfy the 5\% constraint at \emph{any} threshold. We claim no advantage for the judge: its
interval, 0.104 to 0.215, overlaps the mechanical families, and a paired bootstrap of the
conservative-minus-eager difference gives 95\% CI $[-0.058, +0.203]$, spanning zero. What is
precisely measured is the input, not the ranking. The pattern worth a larger study is that whether a
judge is usable at a strict operating point depends on its calibration and not
on its being a judge. An out-of-the-box judge cannot be assumed to respect a strict-budget
stop.

\section{Preregistered Transfer Replication}
\label{sec:transfer}

\begin{figure*}[t]
\centering
\includegraphics[width=0.88\textwidth]{figures/transfer.pdf}
\caption{Confirmatory test on four held-out configurations: no detector and no threshold was fitted
on them before preregistration (they appear in earlier work of ours only through oracle and
run-length quantities, which involve no detector).
Coverage of failed runs against the false-abort budget, per detector family, on a log budget axis.
The shaded band is the deployable 5\% budget. The shape replicates 4 of 4: coverage is near the
floor where the constraint binds, and for the mechanical families it climbs to near-total once the
constraint is lifted. Oracle separation AUC for each configuration is in the panel title.}
\Description{Four side-by-side line plots, one per held-out configuration, showing coverage of
failed runs on the vertical axis against a logarithmic false-abort budget axis. In every panel the
detector curves start near zero coverage at the five percent budget; the mechanical families climb
towards one at the unconstrained budget while self-report often does not.}
\label{fig:transfer}
\end{figure*}

\begin{table*}[t]
\centering
\small
\caption{Every held-out configuration and every detector family, as committed in the frozen
specification. Coverage is the fraction of failed runs on which the family fires at the stated
false-abort budget. ``Transfer'' applies the operating threshold fitted on the development corpus
unchanged; a false-abort rate above 0.05 means the imported threshold violates the budget in
deployment. Rows are reported whether or not they favor the hypothesis.}
\begin{tabular}{llccccccc}
\toprule
& & \multicolumn{4}{c}{Coverage at false-abort budget} & Saved & \multicolumn{2}{c}{Threshold transfer} \\
\cmidrule(lr){3-6}\cmidrule(lr){8-9}
Configuration & Family & 5\% & 10\% & 20\% & 100\% & @5\% & false-abort & coverage \\
\midrule
Claude 3.5 Sonnet, SWE-agent, Verified & Syntactic repetition & 0.049 & 0.049 & 0.049 & 1.000 & 0.008 & 0.000 & 0.012 \\
(82 failed / 47 resolved, AUC 0.684) & Embedding diversity & 0.012 & 0.037 & 0.073 & 1.000 & 0.002 & 0.000 & 0.000 \\
 & Process-reward proxy & 0.171 & 0.244 & 0.293 & 0.988 & 0.086 & 0.021 & 0.134 \\
 & Self-report & 0.000 & 0.085 & 0.085 & 0.976 & 0.000 & 0.638 & 0.780 \\
\midrule
GPT-4o, SWE-agent, Verified & Syntactic repetition & 0.000 & 0.000 & 0.333 & 1.000 & 0.000 & 0.167 & 0.333 \\
(78 failed / 42 resolved, AUC 0.703) & Embedding diversity & 0.000 & 0.000 & 0.321 & 1.000 & 0.000 & 0.190 & 0.321 \\
 & Process-reward proxy & 0.000 & 0.000 & 0.000 & 1.000 & 0.000 & 0.262 & 0.333 \\
 & Self-report & 0.051 & 0.103 & 0.218 & 0.295 & 0.020 & 0.095 & 0.103 \\
\midrule
Claude 4 Sonnet, SWE-agent 2025, Verified & Syntactic repetition & 0.045 & 0.205 & 0.205 & 1.000 & 0.025 & 0.086 & 0.205 \\
(88 failed / 70 resolved, AUC 0.726) & Embedding diversity & 0.068 & 0.068 & 0.170 & 0.920 & 0.061 & 0.014 & 0.011 \\
 & Process-reward proxy & 0.057 & 0.136 & 0.250 & 1.000 & 0.031 & 0.000 & 0.023 \\
 & Self-report & 0.023 & 0.125 & 0.250 & 0.591 & 0.007 & 0.071 & 0.125 \\
\midrule
Qwen3-Coder-480B, OpenHands, SWE-rebench & Syntactic repetition & 0.041 & 0.153 & 0.153 & 1.000 & 0.033 & 0.086 & 0.153 \\
(98 failed / 70 resolved, AUC 0.664) & Embedding diversity & 0.000 & 0.000 & 0.000 & 0.837 & 0.000 & 0.214 & 0.245 \\
 & Process-reward proxy & 0.092 & 0.133 & 0.337 & 1.000 & 0.068 & 0.000 & 0.020 \\
 & Self-report & 0.204 & 0.204 & 0.204 & 0.500 & 0.062 & 0.014 & 0.041 \\
\bottomrule
\end{tabular}
\label{tab:transfer}
\end{table*}

The four held-out configurations were fixed in the frozen specification for diversity of model
family, scaffold and task set: three frontier SWE-agent submissions to the public SWE-bench
experiments leaderboard (GPT-4o and Claude 3.5 Sonnet on the classic tool interface, Claude 4 Sonnet
on the 2025 \texttt{str\_replace\_\allowbreak editor} interface), and Qwen3-Coder-480B under the
OpenHands scaffold \cite{wang2025openhands} on SWE-rebench. All four yielded usable data and all
four are reported. No agent
was run for any of them; all are replay of already published trajectories.

\paragraph{The pattern replicates on all four configurations.} Across the 16 held-out
detector-configuration pairs, all reported in Table~\ref{tab:transfer}, median coverage at the
5\% budget is 0.043 and the maximum is 0.204,
against 0.16 to 0.18 for the mechanical families on the development corpus; five of the sixteen
never fire at all, which is what pulls the median down. Coverage is lower
off-corpus than for the same family on the development corpus in 15 of the 16 pairs, the exception
being self-report on OpenHands. We attach no test to that: the 16 pairs share four corpora and four
detectors and are not independent, the held-out configurations have 42 to 70 resolved runs each so a
5\% budget admits only two to four false aborts, and the corpora differ in composition as well as in
configuration. The defensible statement is that the collapse recurs in every configuration we looked
at, and that the held-out median and 15 of the 16 matched pairs are lower than on the corpus where it
was found. Computing an instance-clustered interval beside each held-out coverage figure sharpens
the collapse rather than softening it. The upper end of all 16 intervals
stays below 0.31 and four are pinned at exactly zero, so even the most optimistic reading of every
held-out pair leaves two thirds of doomed runs uncovered. These four corpora carry at most 1.01 runs
per task instance, so the clustering is inert here and the interval coincides with an ordinary
bootstrap; it does work on the development and matched corpora, at 8.9 and 16 runs per instance. What the intervals do not support is the
finer comparison against the development corpus, where 9 of the 16 lie entirely below the
mechanical band of 0.16 to 0.18 and 7 overlap it, so that ordering stays descriptive. Unconstrained coverage reaches at least 0.92 in 11 of the 12
mechanical pairs (the exception, embedding diversity on OpenHands, reaches 0.837), so the detectors
can see the behavior; what they cannot do is see it only on the runs that are doomed. The rise is
monotone in the budget on every configuration (Figure~\ref{fig:transfer}), so the finding is not
tied to the 5\% operating point, which the frozen specification named as a way it could have failed.

\paragraph{The latency half replicates too.} Coverage alone would leave open whether the held-out
detectors are merely late. They are not. At the 5\% budget, the fraction of failed runs whose alarm
is active at the Oracle Stop is between 0.000 and 0.031 across the 11 pairs that fire at all, and is
exactly zero in 7 of them. Median regret off-corpus is not worth reporting: with between 1 and 20
fired runs per pair it ranges from $-45.5$ to $+57.5$ steps and is dominated by which handful of
runs happened to trip the threshold, itself a symptom of the coverage collapse. We therefore do
not read it as a measurement of latency. Both preregistered hypotheses are supported on all four
configurations. Note that this paragraph, unlike the coverage results, depends on the oracle.

\paragraph{Thresholds do not transfer.} Table~\ref{tab:transfer} also reports the preregistered
threshold-transfer test: each family's operating threshold, fitted on the development corpus, is
applied unchanged to each held-out configuration. Of the 16 pairs, 7 keep the false-abort rate
within the 5\% budget, and the best of those covers 0.134 of failed runs. The other 9 exceed the
budget: among the mechanical families by more than five times (the process-reward proxy on GPT-4o
aborts 26.2\% of resolved runs), and in the worst case overall by nearly thirteen times (self-report
on Claude 3.5 Sonnet aborts 63.8\%).
Across these four case studies a threshold behaves as a property of the deployment and not of
the detector, and on that evidence any reported operating point should be assumed invalid on a new
model, scaffold or task distribution until recalibrated. It is a practical consequence for anyone
shipping a stop-signal, independent of everything the oracle contributes.

\paragraph{What this does not establish.} These are four case studies, not a designed experiment.
Model, scaffold and task set vary together across rows, so nothing here attributes a magnitude to
model capability instead of to a tool interface or a task distribution. The oracle also
discriminates less well on the OpenHands configuration (AUC 0.664 against 0.684 to 0.726
elsewhere), so quantities that depend on $t^{*}$ are localized less confidently there. The coverage
and false-abort columns, which carry the replication, do not depend on $t^{*}$.

\section{Mechanism: Why Coverage Collapses}
\label{sec:mechanism}

\begin{figure*}[t]
\centering
\includegraphics[width=0.88\textwidth]{figures/mechanism.pdf}
\caption{The mechanism, measured without reference to any operating point.
(a)~Fraction of runs whose peak alarm score reaches a given level, for failed and for resolved
runs. The threshold that holds false aborts to 5\% sits deep in a region both populations occupy,
which is why moving it cannot buy coverage.
(b)~Separability of failed from resolved runs, as AUC of the running score, against how far the run
has progressed. It rises to roughly 0.62 by the first third and then stops improving, so waiting
does not sharpen the decision.}
\Description{Two panels. The left panel shows two survival curves, for failed and for resolved
runs, that lie close together across the whole score range, with a dashed vertical line marking the
threshold that holds false aborts at five percent. The right panel shows separability AUC against
fraction of the run elapsed, rising early and then flattening between 0.60 and 0.65 for the
mechanical detectors and near 0.55 for self-report.}
\label{fig:mechanism}
\end{figure*}

A threshold argument alone cannot establish the diagnosis, because coverage and latency are both
readings of one thresholded score. We therefore measure the quantity that any thresholding inherits:
how far apart the two populations are.

\paragraph{Why the run-level grain is the right one.} A detector of the class we evaluate fires on a
run exactly when that run's score first crosses $\theta$, so it fires \emph{at all} precisely when
the run's maximum score reaches $\theta$. Coverage and false-abort at a fixed threshold are
therefore functions of the run-level maximum alone: the run-level analysis below is the frontier's own
quantity with the time index marginalized out. This is
why run-level separation governs the achievable (coverage, false-abort) pairs of any fixed-threshold
rule on the same per-step score, and it is also why the time-resolved curve in
Figure~\ref{fig:mechanism}b, which is the running maximum truncated at each decile, ends exactly at
the run-level AUC in Figure~\ref{fig:mechanism}a. What the argument does not cover is a rule that
reads the shape of the score history, not merely whether it ever crossed a level.

\paragraph{The distributions overlap.} For each run take the maximum alarm score the detector ever
reaches, giving one number per run, and compare the failed and resolved populations
(Figure~\ref{fig:mechanism}a). The two populations are largely the same population to these
signals: the overlap coefficient of the two distributions is 0.75 for
embedding diversity, 0.83 for syntactic repetition, 0.85 for the process-reward proxy and 0.92 for
self-report. Run-level separation is correspondingly weak, at 0.605 (95\% CI 0.536 to 0.669), 0.603
(0.514 to 0.689), 0.617
(0.526 to 0.707) and 0.546 (0.504 to 0.593) AUC respectively, with intervals bootstrapped over task
instances. These are threshold-free statistics. They say that a doomed run and a recoverable run in
this regime look alike to these signals, and that any fixed-threshold rule built on them inherits
that. Setting the threshold at the resolved population's 95th percentile puts 0.16 to 0.23 of failed
runs above it, at a realized false-abort rate of 0.052 to 0.086 (the scores are discrete, so an
exact 5\% rate is not attainable). The frontier's slightly lower operating-point coverages in
Table~\ref{tab:frontier} are the same signal read at thresholds that do respect the 5\% budget.

\paragraph{Waiting does not help.} A natural hope is that the populations separate later, so a
patient detector could do better. They do not (Figure~\ref{fig:mechanism}b). Computing the same
separability at each decile of a run's length, the mechanical families climb from about 0.56 at 10\%
elapsed to a peak of 0.62 to 0.65 between 30\% and 50\%, then flatten or drift down, ending at 0.60
to 0.62. Self-report never exceeds 0.55. There is no late regime where the signal becomes decisive;
by the time a patient detector would have collected more evidence, the evidence has stopped
accumulating and the budget has been spent. This also explains the ten-step Oracle Regret without
recourse to latency: a detector that waits is not waiting for something.

\paragraph{Is the separation just task difficulty?} A pooled comparison mixes two things a deployed
detector must not confuse: that \emph{this run} has stalled, and that \emph{this task} is hard, and
only the first is usable since a detector does not choose its tasks. Our development corpus cannot
separate them, having filled a failed quota and a resolved quota independently, which leaves 7 of
its 80 instances carrying both outcomes.

The public corpus does hold the matched data; the sampler never looked for it. Keeping only
instances with runs of both outcomes yields 348 of them in the first half of the corpus, and from
these we build a \emph{difficulty-matched} corpus: the 60 instances supporting the most
within-instance comparisons, capped at 8 runs per outcome so no task dominates. This gives 960
trajectories (480 failed, 480 resolved) and 3{,}840 within-instance pairs against the development
corpus's 576. Ranking instances by the comparisons they support selects for balanced outcomes, so
tasks the agent always fails and tasks it always solves are absent by construction and every claim
below is bounded to tasks of intermediate difficulty. Selection was fixed before any AUC was
computed and no agent was run.

The confound does not appear (Table~\ref{tab:within}). Holding the task fixed leaves separation
unchanged for every family, by at most $+0.021$ and by $+0.000$ for self-report, and each family's
pooled value falls inside its
within-instance interval. In this matched sample, task difficulty is not what produces the
separation; the earlier 7-instance intervals contained the pooled values throughout, so they never
evidenced a difference. Per-family pooled AUCs do differ between the two corpora (the process-reward
proxy reads 0.617 on the development corpus and 0.572 here), and the best family is not the same one
in both. What the design target rests on is the best value available, and that is nearly unchanged:
0.617 for the process-reward proxy on the development corpus against 0.613 for syntactic repetition
here, so the target's baseline does not move. This analysis is post-hoc, prompted by the confound, and not preregistered
like Section~\ref{sec:transfer}.

\begin{table}[t]
\centering
\small
\caption{Separation holds when the task is held fixed. Measured on a difficulty-matched corpus of 60
instances that each contribute both outcomes (960 runs, 3{,}840 within-instance pairs). Intervals
are bootstrap 95\% CIs over instances. Pooled falls inside the within-instance interval for all
four families, so in this matched sample task difficulty is not what produces the separation.}
\begin{tabular}{lccc}
\toprule
Family & Pooled & Within-instance (95\% CI) & $\Delta$ \\
\midrule
Syntactic repetition & 0.613 & 0.615 \,[0.574, 0.654] & $+0.002$ \\
Embedding diversity  & 0.580 & 0.588 \,[0.549, 0.624] & $+0.008$ \\
Process-reward proxy & 0.572 & 0.593 \,[0.543, 0.641] & $+0.021$ \\
Self-report          & 0.521 & 0.521 \,[0.498, 0.543] & $+0.000$ \\
\bottomrule
\end{tabular}
\label{tab:within}
\end{table}

\paragraph{The claim, precisely.} What the mechanism supports is that these families discriminate
doomed from recoverable long-horizon runs at roughly 0.60 AUC, that task difficulty is not what
produces that separation in a matched sample, and that the level is reached early. It does not establish that no signal could do better, that latency
never matters, or that an online rule reading the shape of the score history could not improve on a
fixed threshold. What it does establish is narrower and still useful: for fixed-threshold rules on
these scores, waiting for a larger value of the same score does not fix low coverage, because the
score stops becoming more discriminative long before the run ends.
\subsection{A Falsifiable Design Target}
\label{sec:target}

The overlap statistic turns the negative result into a specification. Model the run-level alarm
score as Gaussian with unit variance in both populations, separated by $d'$. Setting the threshold
at the resolved population's $(1-\alpha)$ quantile holds false aborts at $\alpha$ and yields
coverage $C = \Phi(d' - z_{1-\alpha})$, so the separation a detector needs is
\begin{equation}
d' = z_{1-\alpha} + z_{C}, \qquad \mathrm{AUC} = \Phi\!\left(d'/\sqrt{2}\right).
\end{equation}
Evaluated at $\alpha = 0.05$, covering a quarter of doomed runs requires
0.754 AUC; half requires 0.878; three quarters requires 0.949. The best family measured here reaches
0.617 (95\% CI 0.526 to 0.707, clustered by task instance), so closing even the gap to half coverage
means adding roughly 0.26 AUC of run-level separation, which is a different signal rather than a
better threshold on this one.

\paragraph{How far can the Gaussian idealization be trusted?} Two checks bound it, and they pull in
opposite directions. First, at the operating points we actually measure the model is pessimistic.
Because the scores are discrete, the threshold at the resolved population's 95th percentile realizes
a false-abort rate of 0.052 to 0.086 instead of exactly 0.05, and at those realized rates every
family covers \emph{more} than the model predicts for its separation, by a factor of 1.4 to 1.7. It
would be tempting to scale the target down by that factor; we do not, because a ratio measured at
coverage near 0.2 licenses nothing about coverage near 0.5.

Second, we redo the derivation without the Gaussian. The threshold is set by the resolved
population, whose distribution we observe directly, so we keep that empirical shape and model an
improved detector as one whose doomed-run scores follow the same distribution shifted up by
$\delta$, the non-parametric counterpart of the equal-variance assumption. Sweeping $\delta$ until
coverage reaches the target gives a requirement of 0.914 to 0.954 across the four families for half
coverage, and 0.79 to 0.93 for a quarter. The empirical tails are heavier than a Gaussian's, so the
threshold sits further out and \emph{more} separation is needed, not less. We therefore quote 0.914
to 0.954 as the requirement and the Gaussian 0.878 as the more forgiving of the two derivations. The two checks disagree about the exact figure by roughly
$0.07$ AUC in either direction and agree completely about the conclusion: every version of the
requirement is far above the 0.617 measured here, and well outside its clustered interval of 0.526
to 0.707.

Two caveats bound this section. The target applies to detectors that fire when a per-step score
first crosses a fixed threshold, the class evaluated here; a rule reading the shape or timing of the
score history is outside this analysis, a route the diagnosis does not close off. And it is
necessary, not sufficient: a signal with the required separation still has to be computable online
and cheaply enough to be worth running. A third limitation, that AUC does not by itself determine
coverage at a given operating point, is taken up in Section~\ref{sec:discussion}.

\section{Does the Instrument Hold Up?}
\label{sec:oracle}
The reference quantity is a proxy, and the strongest objection to this paper is that it measures
lexical overlap with one reference solution rather than progress toward a passing state; we treat
that as a first-class question.

\paragraph{What depends on the oracle, and what does not.} The two results that carry the paper,
the coverage collapse and its 4-of-4 replication, use only the resolved/failed outcome label and the
step at which the detector fires. The mechanism analysis is likewise a comparison of two
outcome-labeled populations; the design target follows from it. The oracle enters in
the Oracle Regret, the PPR reference level and the ``active at $t^{*}$'' statistic. It also enters
in half the support for the comparative claim: showing a detector
was \emph{silent} needs only outcome labels, but showing it was not merely \emph{late} needs a time
reference for what late would mean. The separability plateau supplies that without the oracle; the
``active at $t^{*}$'' statistic beside it does not. A reader who rejects the oracle keeps the
collapse and the plateau and loses the sharper phrasing of the comparison. One
qualification belongs here and not only in the threats: the corpus was filtered with the oracle's
machinery, since a trajectory is kept only when its gold patch yields a non-empty target token set
and at least one of its edit actions can be parsed (the 9 and 29 exclusions of
Section~\ref{sec:design}).
The statistics are oracle-free; the population they are measured on is not. What a reader who
rejects the oracle entirely still faces is detectors that fire on 4\% of doomed runs off-corpus at a
deployable budget.

\paragraph{Calibration.} A valid progress oracle should give resolved runs higher maximum progress
and a later Oracle Stop than failed runs. Resolved runs reach a median gold-token recall of 0.571
against 0.030 for failed runs, separating the classes at 0.762 AUC, and their Oracle Stop lands at
0.42 of the run against 0.21. The two progress views agree on $t^{*}$ for 91.0\% of trajectories;
the 9.0\% contested cases are reported, not silently resolved. Because success is test-defined, a
resolved run can reach its textual peak before the end and still finish by a route the token metric
does not credit, which is why resolved $t^{*}/n$ is 0.42 and not near 1; we therefore make no
progress claim about resolved runs.

\paragraph{The zero-overlap stratum.} Nearly half of failed runs never register any token overlap
with the reference patch, pinning $t^{*}$ at their first edit step, at or near the trajectory
start. The premise that zero overlap
means no progress is testable on runs known to have succeeded: 38 of 291 resolved runs, 13.1\%,
also finish with zero overlap, having solved the task by a route the reference does not describe,
most often by fixing the issue in files the reference patch never touches (34 of the 38).
Zero overlap is thus strongly but not conclusively diagnostic,
$P(\text{failed} \mid \text{zero overlap}) = 0.844$ (Wilson 95\% CI 0.793 to 0.884) against a
corpus base rate of 0.589. This is the error bound a reader should carry into every $t^{*}$-derived
number.

\paragraph{Gold-blind readers land where the token metric does.} The sharpest test of whether
$t^{*}$ tracks progress rather than surface overlap is whether readers who never see the reference
patch agree with it. On 70 failed trajectories across 41 tasks, restricted to the reached-reference
stratum where $t^{*}$ is interior and the judgment is non-trivial, we gave three frontier models
from distinct vendors the issue and the numbered steps, gold patch withheld, and asked each for the
last step of genuine progress. Per judge, the median absolute distance to $t^{*}$ is 2, 1 and 2
steps; the panel median lands 1 step away (n=59 trajectories all three scored), within two steps on
63\% of them, rank-correlated 0.65 to 0.72. This is not a positional prior: guessing the midpoint
misses $t^{*}$ by a median of 5 steps and a uniformly random step by 6.1 to 6.6, so the readers are
three to five times closer than position allows. Clustering by task instance, the mean per-instance
error is 3.0 steps over 38 instances. Where the judges differ from the oracle they place progress
\emph{earlier} in the run, so the token oracle is if anything generous to the agent. We read this as
convergent construct validity, one step short of a second independent ground truth: it is a selected stratum,
the models are the same technology as the agents, and no representative human annotation exists for
the zero-overlap population.

\paragraph{The ``principled'' alternative degenerates.} The cumulative signal is monotone by
construction, so a reviewer will reasonably ask for a current-state oracle whose value can fall when
the agent overwrites correct work. We built one: per gold target file it keeps the agent's active
edit blocks, a new edit drops any block whose line range it replaces, and $t^{*}$ is the last step
attaining the global maximum of line-level recall. On the 121 failed runs it can score, that
alternative places $t^{*}$ at the last applied gold edit in 116 of them, 95.9\%. It reports when the agent last touched a gold file instead of locating a peak,
and that is not a quantity a detector can be scored against. It also validates worse externally. On the 28 trajectories that both oracles
score and all three judges read, the cumulative oracle's median distance to the panel is 3.0 steps
against 5.5, and it is strictly closer on 57.1\% of them against 14.3\% the other way. It is the
natural fix and it fails; the cumulative oracle is the better available instrument on this
evidence.

\paragraph{How often do agents actually regress?} Since the cumulative signal cannot witness an
agent undoing correct work, we measure that separately by edit replay, letting a new edit overwrite
any block whose line range it replaces so alignment is free to fall. One question at two strictnesses, reported as a range: regression occurs in 5.0\% of failed runs at
line level and 23.1\% at identifier level, the latter over the 117 runs with at least two applied
gold edits, with a median largest drop of 0.137 of the reference identifiers, and 16.2\% end below
their own best alignment. Regression is real but a minority: the dominant failure mode is
saturation, an agent that stops improving, not one that reaches a fix and then destroys it, which is
why we read $t^{*}$ as a last-improvement point. Line ranges shift as edits are applied, so the
overwrite model is approximate; we use it to gauge how common regression is, not to define $t^{*}$.

\section{Discussion}
\label{sec:discussion}
\paragraph{What a stop-signal designer should take from this.} The three reporting numbers below are
also three design instructions: report coverage, recalibrate per deployment, and aim at separation
instead of at sensitivity. Only the last carries a target, and for score shapes like those measured
here it is 0.91 to 0.95 run-level AUC to cover half of doomed runs at a 5\% budget. A proposal
that leaves that quantity where it is has to reshape the low false-abort end of the curve to move
the frontier; firing earlier will not do it.

\paragraph{Where the remaining signal might be.} What these families measure, repetition and its
relatives, is common to stuck runs and to hard-but-recoverable ones, so what distinguishes them is
likely to be about the task instead of the surface behavior. Task-grounded signals include how
much of the issue's requirement remains unaddressed, whether the failing test output has changed
in kind, and whether the agent's plan has been revised or merely repeated. Our process-reward stand-in is the crudest version of this and
is already the best or joint-best family at the strict budget, which is weak evidence in that
direction. A trained process reward model \cite{lightman2023verify}, an online changepoint detector
on a task-grounded statistic \cite{adams2007bayesian}, or a judge calibrated to a strict operating
point instead of to accuracy, are the candidates. Testing them needs a fresh confirmatory set: the
four configurations here are spent.

\paragraph{The obvious next experiment, and why it is not here.} All four systems we draw on release
public implementations, so the way to close the gap is to run theirs; the omission is a choice. One
cannot be run on this evidence at all, since log-probability confidence \cite{agentstop2026} reads
token-level log-probabilities that public corpora do not record, and a second needs live execution.
Of the two that are replay-native, the closest fit halts on embedding convergence over successive
drafts in a question-answering loop; carrying that onto agent trajectories means choosing
what counts as a draft and what a plateau means across heterogeneous tool output. Those choices
yield an adaptation, and reporting one under the original's name is the conflation this paper
objects to elsewhere.

\paragraph{The reporting protocol, stated once.} The durable output of this study is the three
numbers a stopping rule should be reported in, each shown load-bearing above and none of them
standard practice today. \textbf{Coverage at a stated false-abort budget}, because savings can look
reasonable while a detector is silent on most doomed runs, and it is the number that made the
collapse visible. \textbf{Run-level separation with an instance-clustered interval}, because it
is what limits coverage on these score shapes and makes results comparable across
corpora of differing difficulty. And \textbf{whether the operating threshold survives a change of
deployment}, because a threshold fitted on one model and scaffold violated the budget on 9 of our 16
held-out pairs, once by nearly thirteen times. None reads the reference patch: the inputs are the
detector's own scores, the resolved/failed labels a test suite already produces, and task
identifiers, plus a second deployment for threshold survival.

\paragraph{What this study does not settle.} The design target is deployment-specific. AUC summarizes an entire ROC curve and does not
determine coverage at a particular low false-abort operating point, which is why the empirical band
sits 0.04 to 0.08 above the Gaussian value. It says what
separation these score shapes require, not what any detector must reach; the protocol therefore asks
for coverage at a stated budget \emph{directly}, with separation beside it, not a substitute.

\section{Threats to Validity}
\emph{The instrument is a proxy.} Distance to the gold patch is not progress toward a passing state,
and the two diverge when a correct fix differs textually from the reference. We restrict the signal
to gold target files and to token recall to reduce this, validate rather than assume (0.762 AUC),
bound the zero-overlap stratum explicitly, publish the contested rate, and show that the natural
current-state alternative degenerates. The residual error is real and bounds the precision of every
$t^{*}$-derived number; it does not touch the coverage, false-abort or run-level separation
results, which use only outcome labels.

\emph{The detectors are operationalizations, not the published systems.} Five families are
reimplemented in spirit, one of them (the process-reward proxy) an explicit stand-in for a trained
model. Conclusions are about these operationalizations as a class of cheap surface-behavior signals,
not about any authors' artifact. A published system with a stronger implementation of the same
family would move the level; it would have to move run-level separation to move the diagnosis. The gap is
closable and we did not close it; Section~\ref{sec:discussion} says which of the four cited
systems could be run against this harness and why we did not.

\emph{No detector ranking is claimed.} With instance-clustered intervals and the threshold
reselected inside each replicate, the three mechanical families are not separated from each other,
and marginal interval overlap is not a pairwise test in either direction. We report the level and
the mechanism, both of which are stable, and no ordering.

\emph{Scope.} Every claim is scoped to long-horizon runs of at least 15 steps, to Python
issue-resolution benchmarks derived from SWE-bench, to the SWE-agent and OpenHands scaffolds, and to
detectors firing when a per-step score first crosses a fixed threshold. Short runs, other languages,
human-in-the-loop settings, and rules reading the shape of a score history instead of its maximum
are untested here.

\emph{Confounded configurations.} The four held-out configurations vary model, scaffold and task set
together, so they establish recurrence, not a capability law. A controlled
replication varying only the model is the priority next step; the harness already ingests both
trajectory formats, so this is data collection, not a method change.

\emph{Truncated runs are not filtered, on either side.} Run-length caps differ across scaffolds and
neither frontier removes runs that hit one, so the comparison is like-for-like rather than
controlled.

\emph{Cost is measured in steps.} The corpus does not expose per-step token counts; we report a
character proxy alongside steps and they agree to within six points.

\emph{The false-abort counterfactual.} We count any alarm before completion on a resolved run as a
lost resolution, overcounting runs that had already produced the passing patch.
Section~\ref{sec:collapse} bounds the effect on the development corpus: exempting alarms after a
resolved run's last applied edit raises the best saved fraction from 0.129 to 0.144, and an
over-generous three-step window raises the best coverage only from 0.18 to 0.24. We did not repeat
this off-corpus, where the resolved sets are smaller (42 to 70 runs) and one reclassified run moves
the budget more, so the transfer results carry the strict convention only.

\emph{The corpus is conditioned on the instrument.} Trajectories are kept only when the gold patch
yields a non-empty target token set and the run contains at least one parseable edit action, so
although coverage and false-abort do not use $t^{*}$, the
population they are measured on was selected with the oracle's machinery. We do not expect this to
bias the resolved-versus-failed comparison, but it is not a fully oracle-free study.

\emph{The gold-blind panel uses paid API models}, unlike the rest of the study, which is free
retrospective replay. It is a one-time validation on 70 trajectories, not part of the main
pipeline; its prompts and raw outputs are in the released artifact.

\emph{Use of AI.} Generative AI tools were used to refine the text; all conceptualization, data
analysis and scientific verification were performed by the authors.

\emph{Data availability.} The harness, frozen specification, cached scores, raw panel outputs
and a verifier that re-derives the paper's numbers are released under an open license at
\url{https://github.com/s1mar/stop-signal-baselines}.

\section{Conclusion}
Stopping rules for long-horizon coding agents are usually described as firing too late. Measured
against a reference-grounded Oracle Stop under a deployable false-abort budget, the four
inexpensive fixed-threshold baselines studied here mostly do not fire at all. The best of them
covers 0.18 of doomed runs on its development corpus, far less on four preregistered held-out
configurations, and its threshold does not survive the move between deployments. Neither the
operating point nor impatience explains this: doomed and recoverable runs score alike at run
level, separability peaks in the first half of a run and never improves, and the separation
survives holding the task fixed in a difficulty-matched sample. This gives the next stopping rule
a number to beat: run-level separation of 0.91 to 0.95 AUC to cover
half of doomed runs at the same budget, for score shapes like those measured here. We would rather
the field take the three reporting numbers than any conclusion of ours about the four baselines we
built: report those, and the collapse is visible in anyone's data. The open problem is not
detecting stagnation. It is telling
stagnation on a doomed run apart from the same behavior on a run that still recovers.
\label{endofcontent}
\bibliographystyle{ACM-Reference-Format}
\bibliography{references}

\end{document}
